\documentclass[12pt,a4paper,twoside]{article}
\usepackage[english]{babel}
\usepackage{amsmath}
\usepackage{draftwatermark}
\usepackage{titling}
\usepackage{graphicx}
\usepackage{import}
\usepackage{subfiles}
\usepackage[section]{placeins}


%evita indentazione a inizio paragrafo
\setlength{\parindent}{0pt}

%fa in modo che le immagini non vengano messe nella subsection successiva
\makeatletter
\AtBeginDocument{%
  \expandafter\renewcommand\expandafter\subsection\expandafter{%
    \expandafter\@fb@secFB\subsection
  }%
}
\makeatother

% Filigrana
\SetWatermarkText{DRAFT}       % Testo della filigrana
\SetWatermarkScale{8}          % Dimensione (scala)
\SetWatermarkColor[gray]{0.9}  % Colore (0.9 è un grigio molto chiaro)
\SetWatermarkAngle{60}       % Angolo (default è 45)

%Title page
\title{AB Inductive Position Sensor}
\author{Simone Vesca}
\date{\today}

\begin{document}

    \begin{titlingpage}
        \begin{center}

            \vspace{2.5cm}

            %Titolo
            \rule{\linewidth}{0.4pt} \\[0.5cm]
            { \huge \bfseries \thetitle }\\[0.4cm] 
            \rule{\linewidth}{2pt} \\[0.8cm]

            %Sottotitolo
            {\large \textit{Design of a Quadrature Encoder, based on 
                            an Eddy Current Inductive Position Sensor}} \\
            %{\large \textit{Progetto di un encoder incrementale AB, basato su 
            % sensore di posizione induttivo a correnti parassite}} \\

            \vfill

            \begin{figure}[h]
                \centering
                \includegraphics[width=0.9\linewidth]{img/biscuit.png}
            \end{figure}

            \vfill

            %Autore E Data
            \begin{minipage}{0.45\textwidth} 
                \begin{flushleft} \large
                    \emph{Autore:}\\
                    \textbf{\theauthor}
                \end{flushleft}
            \end{minipage}
            \begin{minipage}{0.45\textwidth}
                \begin{flushright} \large
                    \emph{Data:}\\
                    \textbf{\thedate}
                \end{flushright}
            \end{minipage}

        \end{center}

    \end{titlingpage}

    %%%%%%%%%%%%%%%%%%%%%%%%%%%%%%%%%%%%%%%%%%%%%%%%%%%%%%%%%%%%%%%%%%%%%%%%%%%%
    % Intro
    %%%%%%%%%%%%%%%%%%%%%%%%%%%%%%%%%%%%%%%%%%%%%%%%%%%%%%%%%%%%%%%%%%%%%%%%%%%%

    \section{Introduction}
    The objective is to design a quadrature incremental position sensor, based on the 
    acquisition of analog signals provided by an {\it IPS (Inductive Position Sensor)},
    and subsequent numerical processing to obtain desired output signals.

    The IPS is an absolute position sensor which outputs voltages proportional 
    to {sine} and {cosine} of the shaft's angular position.
    These signals can be acquired by a microcontroller via an {\it ADC} and 
    processed to extract the position with {\it atan2} operation.

    Finally, the state of the outputs can be obtained dividing the complete 
    revolution into {\it N} angular sectors, corresponding to the desired 
    output {\it PPR (Pulses Per Revolution)}.
    
    \newpage
    \section{Working Principle}
    \subfile{working_principle.tex}

    %%%%%%%%%%%%%%%%%%%%%%%%%%%%%%%%%%%%%%%%%%%%%%%%%%%%%%%%%%%%%%%%%%%%%%%%%%%%
    % Constraints
    %%%%%%%%%%%%%%%%%%%%%%%%%%%%%%%%%%%%%%%%%%%%%%%%%%%%%%%%%%%%%%%%%%%%%%%%%%%%
    \newpage
    \section{Design Constraints}
    %\subfile{working_principle.tex}

    %%%%%%%%%%%%%%%%%%%%%%%%%%%%%%%%%%%%%%%%%%%%%%%%%%%%%%%%%%%%%%%%%%%%%%%%%%%%
    % HW
    %%%%%%%%%%%%%%%%%%%%%%%%%%%%%%%%%%%%%%%%%%%%%%%%%%%%%%%%%%%%%%%%%%%%%%%%%%%%
    \newpage
    \section{Hardware Architecture}
    \subfile{hw_arch.tex}

    \subsection{Coil Design}
    \subfile{hw_arch_coil_design.tex}

    %%%%%%%%%%%%%%%%%%%%%%%%%%%%%%%%%%%%%%%%%%%%%%%%%%%%%%%%%%%%%%%%%%%%%%%%%%%%
    % FW
    %%%%%%%%%%%%%%%%%%%%%%%%%%%%%%%%%%%%%%%%%%%%%%%%%%%%%%%%%%%%%%%%%%%%%%%%%%%%
    \newpage
    \section{Firmware Architecture}
    \subfile{fw_arch.tex}

    \subsection{Normal Mode}
    \subfile{fw_arch_norm_mode.tex}

    \subsection{Calibration Mode}
    \subfile{fw_arch_cal_mode.tex}

    %%%%%%%%%%%%%%%%%%%%%%%%%%%%%%%%%%%%%%%%%%%%%%%%%%%%%%%%%%%%%%%%%%%%%%%%%%%%
    % Conclusions
    %%%%%%%%%%%%%%%%%%%%%%%%%%%%%%%%%%%%%%%%%%%%%%%%%%%%%%%%%%%%%%%%%%%%%%%%%%%%
    \section{Conclusions}

\end{document}